\documentclass[11pt]{article}
\usepackage[margin=1in]{geometry}
\usepackage{amsmath}
\usepackage{booktabs}
\usepackage{array}
\usepackage{longtable}
\usepackage{hyperref}
\begin{document}

\section*{Parameter-to-Equation Map: Density and Holling-II Support}
\textbf{Scope.} This note maps (i) model parameters/equations in
\texttt{Technical\_Reconstruction\_Batch\_DEB\_IBM\_Eisenia\_fetida} and
\texttt{Batch\_DEB\_IBM\_Eisenia\_fetida.m} to (ii) the specific external
papers you asked about.

\subsection*{A) Core model equations used}
\begin{align}
\rho_w &= \frac{N_{adult}+N_{juv}}{\max(S+P_{wet},1)} \tag{60}\\
\rho_{idx} &= \frac{\rho_w}{\rho_w+\rho_{half}} \tag{60}\\
K_S^{eff} &= K_S\left(1+k_{\rho f}\rho_{idx}\right), \quad
f=\frac{S}{S+K_S^{eff}} \tag{61}\\
\frac{dL_{juv}}{dt} &= r_B\!\left(\max(L_i^{common},1.01L_p)-L_{juv}\right) \tag{63}\\
\frac{dL_{adult}}{dt} &= r_B\!\left(L_i^{common}-L_{adult}\right) \tag{64}\\
R_i &= c_R\,\mathrm{reprod\_rate}(L_i,f,p_R) \tag{69}
\end{align}

The chemistry layer is additionally modified by density/activity through
\(\rho_{idx}\) and \(a_{idx}\) (Section 6.12--6.13; Eq. 41--53 in the same reconstruction).

\subsection*{B) Evidence map for density calculations}
\begin{longtable}{p{0.22\textwidth}p{0.34\textwidth}p{0.34\textwidth}}
\toprule
\textbf{Source} & \textbf{What it supports} & \textbf{Mapped model terms/equations} \\
\midrule
\endfirsthead
\toprule
\textbf{Source} & \textbf{What it supports} & \textbf{Mapped model terms/equations} \\
\midrule
\endhead
Yadav and Garg (2016), PMID: 27023819, DOI: 10.1007/s11356-016-6522-7 &
Direct density effect on biological outcomes:
higher stocking density \(\rightarrow\) lower worm biomass/growth/cocoon production; lower density favors biomass growth; mineralization response also changes with density. &
\(\rho_w,\rho_{idx}\) construction (Eq. 60) and density penalty on feeding access through \(K_S^{eff}\) (Eq. 61), which propagates to growth (Eq. 63--64) and reproduction (Eq. 69). \\

Scientific Reports (2022), DOI: 10.1038/s41598-022-17855-z &
Direct density effect on decomposition/mineralization signal, with explicit time-density interaction in their meta-regression. &
Supports keeping density as an explicit predictor in process-level soil/organic-matter behavior; mapped to \(\rho_{idx}\)-driven chemistry modifiers in Section 6.12--6.13 (Eq. 41--53). \\
\bottomrule
\end{longtable}

\subsection*{C) Evidence map for Holling Type II functional response}
\begin{longtable}{p{0.22\textwidth}p{0.34\textwidth}p{0.34\textwidth}}
\toprule
\textbf{Source} & \textbf{What it supports} & \textbf{Mapped model terms/equations} \\
\midrule
\endfirsthead
\toprule
\textbf{Source} & \textbf{What it supports} & \textbf{Mapped model terms/equations} \\
\midrule
\endhead
van der Meer (2016), Conservation Physiology 4(1):cow023, DOI: 10.1093/conphys/cow023 &
Direct DEB-form functional-response framing: scaled functional response \(f\) as the food-density linkage and critical resource-density logic in DEB population context. &
Justifies using the saturating food-response structure
\(f=S/(S+K_S^{eff})\) in Eq. (61), and its downstream use in growth/reproduction Eq. (63--64, 69). \\

Yadav and Garg (2016), PMID: 27023819 &
Indirect support only for Holling-II form: shows density-linked performance shifts consistent with resource competition, but does not derive Holling-II mathematically. &
Use as empirical plausibility for competition/crowding effects; do \emph{not} cite as primary derivation of Eq. (61). \\
\bottomrule
\end{longtable}

\subsection*{D) Recommended citation-strength statement}
\begin{itemize}
\item \textbf{Primary/formal support for Holling-II form}: DEB/conservation physiology framing (cow023) plus DEB literature.
\item \textbf{Primary empirical support for density term calibration need}: PMID 27023819 and 10.1038/s41598-022-17855-z.
\item \textbf{Most defensible wording in report}: ``Holling-II response is adopted from DEB-style feeding formulation; density modifier is empirically motivated by earthworm stocking-density and density-mineralization studies.''
\end{itemize}

\subsection*{E) Quick map of the three constants from your code snippet}
\begin{align}
R_i &= c_R\,\mathrm{reprod\_rate}(L_i,f,p_R), \quad c_R=\texttt{reprod\_correction} \tag{69}\\
i_{hatch} &= \left\lfloor \frac{a_b^T}{\Delta t}\right\rfloor+1, \quad a_b^T=\texttt{egg\_incubation\_days} \tag{28}\\
N_{new} &\sim \mathrm{Binomial}\!\left(\mathrm{round}(E_{hatch}^{inc}),\,s_{hatch}\right), \quad s_{hatch}=\texttt{hatch\_surv} \tag{29}
\end{align}
These additionally enter the horizon egg-control expectations \(E[H_{inc}],E[H_{bin,all}]\) and therefore Eq. (55)--(57), compendium Eq. (80)--(82).

\subsection*{Links}
\begin{itemize}
\item \url{https://pubmed.ncbi.nlm.nih.gov/27023819/}
\item \url{https://www.nature.com/articles/s41598-022-17855-z}
\item \url{https://academic.oup.com/conphys/article/4/1/cow023/2951363}
\end{itemize}

\end{document}
